%==============================================================================
% Section 4 -- Colour Modelling with Spherical Harmonics
%==============================================================================
\section{Colour Modelling via Spherical Harmonics}
\label{sec:sh}

\subsection{Motivation}

Surfaces exhibit view-dependent appearance (specular highlights, metallic
reflections). In 3DGS each Gaussian encodes its view-dependent colour compactly
using \textbf{spherical harmonics (SH)}, an orthonormal basis on the unit
sphere $S^2$

%------------------------------------------------------------------------------
\subsection{View-Dependent Colour}

For Gaussian $\Gcal_i$ and channel $k \in \{R,G,B\}$, the colour in direction
$\mathbf{d}$ is:
\begin{equation}
  \boxed{
    c_{i,k}(\mathbf{d}) = \sum_{l=0}^{L} \sum_{m=-l}^{l}
      f_{i,k,l}^m\, Y_l^m(\mathbf{d})
  }
  \label{eq:sh_color}
\end{equation}
where $f_{i,k,l}^m$ are learned coefficients. In matrix form:
$\mathbf{c}_i(\mathbf{d}) = \mathbf{F}_i\,\mathbf{y}(\mathbf{d})$,
with $\mathbf{y}(\mathbf{d}) \in \R^{(L+1)^2}$ the vector of SH values.
SH evaluation is purely polynomial, making it very efficient on the GPU.

%------------------------------------------------------------------------------
\subsection{Degree Selection}

\begin{table}[H]
  \centering
  \caption{SH parameter count by degree.}
  \label{tab:sh_count}
  \renewcommand{\arraystretch}{1.3}
  \begin{tabular}{cccc}
    \toprule
    \textbf{Max Degree $L$} & \textbf{Coeffs/channel} & \textbf{Total (RGB)} & \textbf{Effect} \\
    \midrule
    0 & 1  & 3  & Constant diffuse \\
    1 & 4  & 12 & Linear shading \\
    2 & 9  & 27 & Quadratic specular \\
    3 & 16 & 48 & Sharp specular (used in 3DGS) \\
    \bottomrule
  \end{tabular}
\end{table}

The original 3DGS paper uses $L=3$. Training starts at $L=0$ and unlocks one
degree every 1000 iterations, stabilising early optimisation by fitting
low-frequency geometry first.

\newpage
